% ---------------------------------------------------------------------
% Preguntas (lo editan los tres)
% 1) Preguntas que esperan del docente o de otros equipos: cada una con
%    quién responde primero, aunque cualquiera debe poder hacerlo.
%    Practíquenlas cruzadas (Erwin responde de los paneles, César de TCP...).
% 2) La pregunta crítica que el grupo le hará a otro equipo (cuenta en D5).
% ---------------------------------------------------------------------
% xltabular: la tabla puede partirse entre páginas y repite el encabezado.
{\small\renewcommand{\arraystretch}{1.2}
\begin{xltabular}{\textwidth}{@{}>{\raggedright\arraybackslash}p{6.0cm}>{\raggedright\arraybackslash}p{2.4cm}X@{}}
  \caption{Preguntas previstas}\label{tab:preguntas}\\
  \toprule
  \textbf{Pregunta probable} & \textbf{Responde primero} & \textbf{Idea de la respuesta}\\
  \midrule
  \endfirsthead
  \toprule
  \textbf{Pregunta probable} & \textbf{Responde primero} & \textbf{Idea de la respuesta}\\
  \midrule
  \endhead
  \bottomrule
  \endlastfoot
  \pend{Grupo}{} & & \\
  \pend{Grupo}{} & & \\
  \pend{Grupo}{} & & \\
\end{xltabular}}

\noindent\textbf{Pregunta crítica para otro equipo:} \pend{Grupo}{a qué equipo y qué le van a preguntar}
